% ==========================================================================
% GENERATED by scripts/split_paper.py. Do not edit this file.
%
% Source: imported_papers/grokking/iclr2027_conference_v2.tex
% Part 5: Relevance: Distributed circuits tolerate neuron deletion, concentrated circuits keep spare capacity
%
% Edits belong upstream, in the published paper. To change what the
% chapter includes, change imports/grokking.toml and re-run
% the splitter.
% ==========================================================================
\section{Relevance: Distributed circuits tolerate neuron deletion, concentrated circuits keep spare capacity}
\label{sec:fragile}

Throughout this section the concentrated circuit is the reference ($\beta = 0$) and the distributed circuit is the fully equalized network ($\beta = 1$), both trained on modular addition at $p = 31$.


\subsection{Tolerance to neuron deletion}
A computation spread over many neurons should depend less on any single one of them \citep{georgopoulos1986population}, and networks that generalize well are known to rely less on single units \citep{morcos2018importance}. In an additive model, the relative fluctuation of the output under random deletion falls with the effective neuron count (Proposition~\ref{prop:fluct}, Appendix~\ref{app:theory}). We test this with three deletion procedures, all applied to the first hidden layer of the trained network without retraining. Random deletion removes a random subset of the neurons. Norm-ranked deletion removes neurons in decreasing order of outgoing norm, so the neurons a concentrated circuit relies on go first. Greedy deletion removes, at each step, the neuron whose removal lowers test accuracy the most. For random and norm-ranked deletion we report the \emph{neuron-deletion tolerance}, or \emph{deletion tolerance} for short, the normalized area under the curve of test accuracy against the deleted fraction from $0$ to $0.8$. For greedy deletion we report the number of neurons removed before test accuracy falls below $0.5$ (Appendix~\ref{app:deletion}).

\begin{table}[t]
\caption{Deletion tolerance and adaptation of the concentrated and the distributed circuit on modular addition ($p=31$). Means over three seeds for deletion, five for pruning to test accuracy $0.9$ (Appendix~\ref{app:deletion}), ten for subtraction and five for multiplication, with standard deviations in parentheses (Table~\ref{tab:adapt}).}
\label{tab:consequences}
\begin{center}
\footnotesize
\begin{tabular}{lccc}
\toprule
 & Concentrated ($\beta=0$) & Distributed ($\beta=1$) & From scratch \\
\midrule
Random deletion tolerance & $0.674$ & $0.770$ & --- \\
Norm-ranked deletion tolerance & $0.332$ & $0.581$ & --- \\
Greedy deletion, neurons removed & $23$ & $61$ & --- \\
Magnitude pruning, fraction removed & $0.62$ & $0.20$ & --- \\
Greedy pruning, neurons removed & $191$ & $160$ & --- \\
Adaptation to subtraction, epochs & $640\,(128)$ & $1540\,(180)$ & $1410\,(137)$ \\
Adaptation to multiplication, epochs & $1800\,(141)$ & $1940\,(136)$ & $1200\,(63)$ \\
\bottomrule
\end{tabular}
\end{center}
\end{table}

The distributed circuit is more tolerant under all three procedures (Table~\ref{tab:consequences}). The advantage is largest under norm-ranked deletion, as expected if the concentrated circuit rests on few high-norm neurons, and it persists under greedy deletion. In this network, the tolerance follows the distribution of outgoing norms rather than the rule that produced it: an outgoing-norm constraint that holds participation at $0.99$ without changing the gradients gives the same random-deletion tolerance as equalization (Appendix~\ref{app:suff}). Nero reaches similar deletion tolerances without the intervention (Appendix~\ref{app:zoo}).

Removing neurons in the opposite order, smallest outgoing norm first as in magnitude pruning, favors the concentrated circuit. It keeps test accuracy above $0.9$ until $62\%$ of its neurons are removed, against $20\%$ for the distributed circuit, because its nearly unused neurons can go at no cost. Greedy pruning gives the same order (Table~\ref{tab:consequences}). The circuit that pruning extracts for the same algorithm also depends on training. At width $2048$, where both start at test accuracy $0.998$, greedy pruning to $0.95$ keeps $33$ to $98$ neurons of the reference and $138$ to $229$ under equalization, and magnitude pruning keeps $143$ against $1692$ on average. A decayed per-neuron gain lowers the greedy circuit to $14$ to $22$ neurons, at least $7$ times fewer than under equalization in each seed (Appendix~\ref{app:minimal}).

\subsection{Adaptation to a new operation}
We take networks that have grokked addition at each $\beta$ and fine-tune the whole network on a new operation with backpropagation with Adam. The \emph{adaptation time} is the number of fine-tuning epochs until test accuracy on the new operation reaches $0.9$. We compare it with training a new network from scratch.

We first transfer to subtraction, which uses the same Fourier frequencies as addition and differs in how they are combined. Training from scratch takes $1410$ epochs. Starting from the concentrated addition circuit takes $640$ epochs, so the pretrained circuit helps. Starting from the distributed circuit takes $1540$ epochs, no faster than from scratch, and the adaptation time increases monotonically with $\beta$ in between (Appendix~\ref{app:adapt}). The concentrated circuit adapts about $2.4$ times as fast as the distributed one. The ordering holds at one third and at three times the fine-tuning learning rate, in the reverse transfer from subtraction to addition, when an outgoing-norm constraint sets the participation of the source network, and when both networks are fine-tuned with credit equalization instead of the reference rule (Appendices~\ref{app:suff} and~\ref{app:adapt}). The difference also appears when both continue from the same grokked checkpoint, with $1052$ against $608$ epochs, slower in all five seed pairs.

Much of the advantage comes from the nearly unused neurons. With them removed from the concentrated network at unchanged test accuracy, adaptation to subtraction takes $1204$ instead of $652$ epochs, against $1480$ for the distributed network, so they account for about two thirds of the difference. On $(b - a) \bmod p$ they account for about half (five seeds, Appendix~\ref{app:adapt}). Networks distributed by per-neuron decay alone adapt as slowly as equalized ones, so across the tested interventions adaptation follows the circuit, not the rule that produced it. Targets that keep the hidden features of addition and change only the readout, $(a + b + 7) \bmod p$ and a fixed permutation of $(a + b) \bmod p$, are learned $2.3$ to $6.6$ times faster than from scratch, with no consistent advantage for either circuit.

We repeat the experiment with multiplication, which does not share the Fourier frequency representation of addition. Every pretrained network takes $1800$ to $2020$ epochs, compared with $1200$ from scratch, so pretraining on addition does not help at any $\beta$.
